

Similar to toric varieties, for a spherical $G$-variety $X$, there are only finitely many $B$-invariant prime divisors, see \cite[Theorem 2.2.2]{Pezzini_lec_sph_wond_vari}. Let
\begin{itemize}
    \item $X_1, \dots, X_k$ be the $G$-invariant prime divisors on $X$, and
    \item $Y_1, \dots, Y_{k'}$ be the $B$-invariant but not $G$-invariant prime divisors on $X$.
\end{itemize}
\begin{remark}
    The prime divisors $Y_1, \dots, Y_{k'}$ are called \emph{colors} of a spherical variety.
\end{remark}

As the $\mathbb{Q}$-divisor $D^B_X$ is $B$-invariant, we can find the coefficients of $D^B_X$ by 
\begin{equation}
    D^B_X = \sum_{i = 1}^{k'} \alpha_i Y_i + \sum_{j = 1}^{k} \beta_j X_j
\end{equation}
where $\alpha_i, \beta_j \in \mathbb{Q}_{\geq 0}$.

Similar to \Cref{cor:tor:Kss}, we have a weaker result for spherical varieties,
\begin{corollary}\label[corollary]{cor:sph:Kss}
	If a $\mathbb{Q}$-Fano spherical variety $X$ is K-semistable, then we have the inequality $\beta_j \leq 1$ for every $j = 1, \dots, k$.
\end{corollary}
\begin{proof}
    This follows directly from \Cref{thm:sph} and \Cref{thm:Zhuang}. For every $j$, the coefficient $\beta_j = \ord_{X_j}(D^B_X)$, and as divisors $X_1, \dots, X_k$ are prime divisors on $X$, we have $A_X(\ord_{X_j}) = 1$.
\end{proof}

There is another natural $B$-invariant representative for anticanonical $\mathbb{Q}$-divisors, which is indeed a $\mathbb{Q}$-Cartier divisor, for a (klt) Fano spherical $G$-variety $X$.
\begin{theorem}[c.f., {\cite[Theorem 1.2]{GH_gor_sph_fano_var}}, based on {\cite[\S4.1, 4.2]{Brion_cur_div_sph_var}}]
    There is a $B$-invariant anticanonical divisor of the form
    \begin{equation*}
        -\cano_X := \sum_{i = 1}^{k'} m_i Y_i + \sum_{j = 1}^{k} X_j,
    \end{equation*}
    where the coefficients $m_j \in \mathbb{Z}_{> 0}$ depend only on the open $G$-orbit of variety $X$.
\end{theorem}
\begin{remark*}
    We call this divisor \emph{Brion's anticanonical divisor}. For toric varieties, Brion's anticanonical divisor is the same as the standard anticanonical divisor $X_1 + \dots + X_k$.
\end{remark*}

Following \Cref{sec:prelim:bir:dis}, we can use Brion's anticanonical divisor to define the log discrepancy.

There is a deep relationship between Brion's anticanonical divisor and the $\mathbb{Q}$-divisor $D^B_X$ for K-semistable varieties.
\begin{corollary}
    Given a $\mathbb{Q}$-Fano spherical $G$-variety $X$, $X$ is K-semistable iff for every valuation $v \in \GDValo_X$, the equality $v(D^B_X) \leq v(-\cano_X)$ holds. Specifically, if variety $X$ is horospherical, which means the valuation cone is the whole space, then $X$ is K-semistable iff $D^B_X = -\cano_X$.
\end{corollary}
\begin{remark*}
    This is the reformulation of Delcroix's criteria (see \cite[Theorem 5.3 and Corollary 5.7]{Delcroix_Kst_Fano_spher}). Note that the anticanonical divisor $-\cano_X \in \Qls{L}^B$ here plays the same role as the point $2 \rho_Q \in \Delta^+$ in \textit{loc.~cit.}
\end{remark*}
\begin{proof}
    According to \cite[Proof of Proposition 5.2]{Pas_sing_sphe}, for every valuation $v \in \GDValo$, we have $v(-\cano_X) = A_X(v)$. The result then follows directly from \Cref{thm:sph}.
\end{proof}

Finally, let us study the coefficient $\alpha_i$ of the $\mathbb{Q}$-divisor $D^B_X$. If the image of valuation $\ord_{Y_i}$ is in the valuation cone, then we have the inequality $\alpha_i \leq m_i$. In general, coefficient $\alpha_i$ may be larger than $m_i$.
\begin{example}
    Let $G := \SL_2(\mathbb{C})$, and the Borel subgroup $B$ is the subgroup of upper triangular matrices. The variety $\mathbb{P}^1 \times \mathbb{P}^1$ is spherical under the diagonal action. The variety $X$ has three $B$-invariant prime divisors,
    \begin{equation*}
    \begin{aligned}
        Y_1 &:= \{[1: 0]\} \times \mathbb{P}^1, \qquad Y_2 := \mathbb{P}^1 \times \{[1: 0]\}, \\
        X_1 &:= \{(x, x) \in \mathbb{P}^1 \times \mathbb{P}^1 \mid x \in \mathbb{P}^1 \}.
    \end{aligned}
    \end{equation*}
    The anticanonical bundle is $\stsh_X(2, 2)$, so that $-\cano_X = Y_1 + Y_2 + X_1$, and according to \cite[Example 5.10]{Delcroix_Kst_Fano_spher}, we have $D^B_X = {4 \over 3} Y_1 + {4 \over 3} Y_2 + {2 \over 3} X_1$.
    
    \begin{figure*}[h]
        \begin{tikzpicture}[
    dot/.style={circle, fill, inner sep=1.5pt}, % 定义点的样式
    label distance=2pt
]
    % 2. 绘制主线段
    \coordinate (A) at (0,0);
    \coordinate (B) at (6,0);
    \draw[thick] (A) -- (B);

    % 3. 标注并绘制点
    % 端点
    \node[dot, label=above:{$2 X_1$}] at (A) {};
    \node[dot, label=above:{$2 Y_1 + 2 Y_2$}] at (B) {};

    % 中点 (1/2)
    \node[dot, label=above:{$-\cano_X$}] (M) at ($(A)!0.5!(B)$) {};

    % 2/3 点
    \node[dot, label=below:{$D^B_X$}] (T) at ($(A)!0.6667!(B)$) {};

\end{tikzpicture}
        \caption{Rational linear series $\Qls{{-\cano_X}}^B$.}\label[figure]{fig:P1xP1}
    \end{figure*}

    In this case, we have the inequality $\alpha_1 = {4 \over 3} > 1 = m_1$.
\end{example}
